% Generated draft: requires mathematical and source review.
\documentclass[11pt]{article}
\usepackage[margin=25mm]{geometry}
\usepackage{fontspec}
\setmainfont{Latin Modern Roman}
\usepackage{amsmath,amssymb,mathtools,unicode-math}
\setmathfont{Latin Modern Math}
\usepackage{enumitem,needspace,xurl,hyperref}
\setlength{\parindent}{0pt}
\setlength{\parskip}{5pt}
\setlength{\emergencystretch}{3em}
\newcommand{\N}{\mathbb{N}}
\newcommand{\Z}{\mathbb{Z}}
\newcommand{\Q}{\mathbb{Q}}
\newcommand{\R}{\mathbb{R}}
\newcommand{\C}{\mathbb{C}}
\newcommand{\F}{\mathbb{F}}
\newcommand{\A}{\mathbb{A}}
\newcommand{\PP}{\mathbb{P}}
\newcommand{\E}{\mathbb{E}}
\newcommand{\eps}{\varepsilon}
\newcommand{\dd}{\,\mathrm{d}}
\begin{document}
\section*{UnsolvedMath ID: 30003497}
\textbf{Source problem code:} catalogue entry
\subsection{Definitions and mathematical statement}
Under the usual interpretation, a real entire function is a holomorphic function \(F\colon\C\to\C\) satisfying \(F(\R)\subseteq\R\). Its associated expansion of the real field is
\[
\mathcal R_F=(\R;0,1,+,\cdot,<,F|_{\R}).
\]
The supplied record does not specify a narrower class of functions or additional hypotheses on \(F\).

If \(F\not\equiv0\), put
\[
U_F=\{z\in\C:F(z)\ne0\},\qquad
G\colon U_F\to\C,\qquad
G(z)=\frac{\mathrm d}{\mathrm dz}\left(z\frac{F'(z)}{F(z)}\right).
\]
Here primes denote complex differentiation. This formula defines \(G\) on \(U_F\); its real restriction is initially defined on \(U_F\cap\R\). The source does not state the necessary exclusion of \(F\equiv0\), specify treatment of zeros of \(F\), or say whether any removable singularities should be filled in. Thus this domain describes where the displayed formula is directly meaningful, rather than an asserted intended domain for the source question.

\textbf{Source question.} For an expansion of the real field by a real entire function \(F\), what technical, mathematical, and model-theoretic advantages arise from studying
\[
G(z)=\left(z\frac{F'(z)}{F(z)}\right)'
\]
instead of \(F\) itself?

\textbf{Unresolved formulation.} The record supplies an explanatory question, without identifying a property, construction, comparison criterion, or conclusion to be established. It does not specify whether studying \(G\) means examining a complex function, its real restriction, or a structure expanded by that restriction. Consequently, a precise quantified conjecture cannot be recovered from the supplied evidence. The definitions above clarify a conventional reading of the displayed expression; they do not supply the missing intended target.
\subsection{Short English statement}
The question asks why the transform \(G=(zF'/F)'\) is useful when studying a real field expanded by a real entire function \(F\). The source does not identify the advantages to prove, the required hypotheses, or the intended domain of \(G\), so this remains a qualified statement of the question requiring human mathematical review.
\subsection{Sources}
\begin{itemize}
\item Supplied local UnsolvedMath registry record, permanent ID \(30003497\), title: \emph{Model-Theoretic Advantages of Logarithmic-Derivative Transforms}; source problem code \texttt{catalogue entry}. Collection: \emph{Oberwolfach Reports Open Problems}; category: Logic. No individual proposer is supplied; the proposed year is recorded as 2017.
\item Supplied report citation: \emph{O-Minimality and its Applications to Number Theory and Analysis} (2018), Oberwolfach Reports, pp.\ 1349--1420, DOI \texttt{10.4171/owr/2017/22}. Supplied source URL and literature-source link, labelled ``Oberwolfach Reports 2018'': \url{https://doi.org/10.4171/owr/2017/22}. The background attributes the question to workshop open-problem discussions documented in Oberwolfach Reports. No question-specific page or speaker attribution is supplied.
\item Literature assessment dated \texttt{2026-08-21}: the supplied status is ``open,'' while the accompanying assessment says that the fragment lacks the relevant construction or claim and therefore permits no definite resolution assessment. This is supplied dated curation, not an independently verified literature assessment. The background expressly requests independent verification of time-sensitive status and cited claims.
\item The registry labels the statement \texttt{corrected\_verified} and reports that mathematical content was preserved while notation, typography, or omitted definitions were repaired. This is supplied editorial metadata, not independent mathematical verification. The report citation uses 2018, whereas the DOI contains 2017 and the proposed year is 2017; their relationship has not been verified here.
\end{itemize}
\end{document}
